\section{Week 9: rANS Byte Stack and Actual Core Boundary}

Week~9 has one target, \claimid{OBL-RANS-CORE-INVERSE}.  The sprint does not
reprove the Week~8 coefficient leaves, table certificate, pure step inverse,
or full-procedure extraction identity.  It fixes a common trace for the
reverse encoder and forward decoder, proves the normalization arithmetic and
the actual suffix copy, and places the three actual generated calls in one
unary harness.  The two generated outer-loop refinements remain open, so the
result is \status{continue-rans}, not a completed inverse.

\subsection*{Pinned actual procedures and hashes}

The proof uses the Week~8 focused targets without a second extraction:

\begin{center}
\begin{tabular}{l l}
\toprule
Generated target & SHA-256 prefix \\
\midrule
\path{RansEncodeTarget.ec} & \texttt{3184367f5d41196b} \\
\path{RansDecodeTarget.ec} & \texttt{2659860ffe3dff9d} \\
\path{HbzFullEncodeTarget.ec} & \texttt{6c085864d31a3cb4} \\
\bottomrule
\end{tabular}
\end{center}

The source hashes remain \texttt{c2caff5b\ldots e1f4e9} for
\path{hpoly.jazz} and \texttt{e800d1d9\ldots b5cbd} for
\path{encoding.jazz}. The generated procedures are:
\begin{itemize}
\item \path{RansEncodeTarget.M._rans_encode};
\item \path{HbzFullEncodeTarget.M.__copy_encoded_suffix};
\item \path{RansDecodeTarget.M._rans_decode}.
\end{itemize}
Extraction regeneration and exact generated hashes are checked before proof
compilation.

\subsection*{Reverse/forward state correspondence}

For symbols $s_0,\ldots,s_{n-1}$, $n=1024$, the boundary-state convention is
\[
\begin{aligned}
  x_n&=2^{23},\\
  x_0&=\text{the decoder's serialized initial state},\\
  \text{state after decoding }k\text{ symbols}&=x_k.
\end{aligned}
\]
The encoder processes $j=n-1,\ldots,0$. At step $j$ it normalizes $x_{j+1}$,
applies the fast concrete-table step, and obtains $x_j$. The decoder processes
$k=0,\ldots,n-1$, selects $s_k$, applies the inverse arithmetic step, and
consumes the byte segment that reconstructs $x_{k+1}$.

Cursor boundaries use
\[
  c_0=4,\qquad
  [c_k,c_{k+1})=\text{normalization bytes for symbol }s_k,\qquad
  c_n=\mathsf{size}.
\]
The first four bytes are $\mathsf{serialize32\_le}(x_0)$. In EasyCrypt these
objects are \path{trace_states}, \path{trace_segments}, \path{trace_cuts},
and \path{trace_bytes}; \path{valid_rans_trace} merely packages their defining
equalities and is not an implementation assumption.

\subsection*{Normalization-byte stack and byte order}

For the mode-2 table,
\[
  2^{23}\leq x<2^{31},\qquad
  x_{\max}(s)=2^{21}f_s\geq2^{21}.
\]
Thus the chosen pure normalization emits zero, one, or two bytes. If
$x=2^{16}q+256h+\ell$, two encoder iterations write $\ell$ first and $h$
second while moving the cursor backward. The final decoder-order segment is
therefore $[h,\ell]$, and
\[
  q\xrightarrow{h}256q+h\xrightarrow{\ell}2^{16}q+256h+\ell=x.
\]

The compiled leaf results are separated as follows.
\begin{itemize}
  \item \path{renorm_bytes_readback} proves the list-level equation;
  \item \path{encoder_shift8_uint} and \path{encoder_low_byte_uint} identify
    actual W32 shift/truncate with division/remainder; and
  \item \path{decoder_append_encoder_byte} proves the one-byte W32 inverse;
    and
  \item \path{encoder_two_byte_decoder_order} proves the two-byte W32 inverse.
\end{itemize}
The W64 cursor lemmas prove decrement safety under their stated bounds.

The pure state invariant is also compiled:
\[
  2^{23}\leq x_{j+1}<2^{31}
  \Longrightarrow
  2^{23}\leq x_j<2^{31}.
\]
It uses the Week~8 concrete reciprocal-step theorem after proving the reduced
state lies in $[1,x_{\max}(s))$.

\subsection*{Encoder and decoder loop invariants}

The intended actual encoder outer invariant at counter $i$ is
\[
\begin{aligned}
  0&\leq i\leq1024,\\
  x&=\mathsf{encode\_trace}(s_i,\ldots,s_{1023}).\mathsf{state},\\
  \mathsf{encp}[off,1024)
    &=\mathsf{traceBytesWithoutState}(s_i,\ldots,s_{1023}),\\
  off+\lvert\mathsf{traceBytesWithoutState}\rvert&=1024.
\end{aligned}
\]
The inner invariant must additionally relate zero, one, or two actual shifts
and backward stores to the corresponding prefix of the pure normalization
segment. On a live iteration it must retain enough cursor space for the final
four-byte state and frame all unwritten bytes.

The intended actual decoder invariant is
\[
\begin{aligned}
  0&\leq i\leq1024,\\
  \mathsf{symsp}[0,i)&=(s_0,\ldots,s_{i-1}),\\
  x&=x_i,\qquad off=c_i,\qquad bad=0,\\
  \mathsf{bufp}[c_i,\mathsf{size})&=\text{the remaining trace suffix}.
\end{aligned}
\]
At exit it should imply $x=2^{23}$, $off=\mathsf{size}$, exact symbol recovery,
decoder success, and the output-array frame. Neither actual outer invariant is
yet compiled.

\subsection*{Actual suffix-copy refinement}

The previous quantified-exit residual is closed. The generated-procedure
theorem is
\begin{center}
  \path{copy_encoded_suffix_correct}.
\end{center}
Its premises are
\[
  0\leq off_0,\quad0\leq n,\quad off_0+n\leq2048
\]
and concludes
\[
\begin{aligned}
  \forall j\in[0,n):\quad
    out_{\mathrm{after}}[j]&=enc_{\mathrm{before}}[off_0+j],\\
  \forall j\in[n,2048):\quad
    out_{\mathrm{after}}[j]&=out_{\mathrm{before}}[j].
\end{aligned}
\]
The proof opens the actual copy loop. It does not construct a desired decoder
buffer. Deriving its bounds from the actual encoder-success branch remains in
the encoder refinement.

\subsection*{Success/failure disjunction and actual harness}

The actual encoder and decoder each have a direct control theorem fixing the
concrete mode-2 tables and parameters and proving the published bad word is in
$\{0,1\}$. These are implementation theorems, but their conclusion is control
only rather than trace refinement.

The unary \path{Mode2RansActualHarness.run} executes
\[
\procname{_rans_encode}
\longrightarrow
\begin{cases}
  bad\neq0 &: \mathsf{decoder\_ran}=\mathsf{false},\\
  bad=0 &: \procname{__copy_encoded_suffix}(off,1024-off)\\
        &\quad\longrightarrow
          \procname{_rans_decode}(1024,1024-off,13).
\end{cases}
\]
The theorem \path{actual_rans_harness_branches_on_encoder_result} proves the
branch equivalence and exact decoder input fields. Encoder success is obtained
from the actual return value rather than assumed. It does not conclude decoder
success or symbol equality.

The desired final disjunction remains
\[
\begin{aligned}
  bad_{enc}\neq0\quad\lor\quad\bigl(&
    decoded[0,1024)=symbols[0,1024)\\
    &\land x_{final}=2^{23}
     \land off_{final}=\mathsf{size}
     \land \mathsf{frames}\bigr).
\end{aligned}
\]

\subsection*{Non-vacuity and exact proof status}

The pure symbol/state premises and the decoder-state construction have
compiled witnesses. No EasyCrypt theorem yet proves that the all-6 symbol
array makes the actual encoder return success with
$4\leq\mathsf{size}\leq1024$. Hence
\claimid{OBL-RANS-ACTUAL-SUCCESS-WITNESS} remains \status{partial}; runtime
execution is not used as a theorem.

\begin{center}
\begin{tabular}{l l l}
\toprule
Obligation & Status & Boundary \\
\midrule
Normalization byte inverse & \status{proved} & pure/word leaf \\
State-bound preservation & \status{proved} & pure trace \\
Actual suffix copy & \status{proved} & generated procedure \\
Actual harness control & \status{proved} & three actual calls, no inverse \\
Actual encoder refinement & \status{partial} & trace postcondition open \\
Actual decoder refinement & \status{partial} & trace consumption open \\
Actual core inverse & \status{partial} & no round-trip conclusion \\
Actual success witness & \status{partial} & no compiled witness \\
\bottomrule
\end{tabular}
\end{center}

\subsection*{Failed tactics and next gate}

A weak outer-loop invariant can prove $bad\in\{0,1\}$ and can preserve every
early/failure branch. It cannot synthesize the array/list byte-stack relation.
Direct encoder/decoder lockstep was rejected because their directions and
inner-loop counts differ. A constructed decoder buffer, external
\path{valid_rans_trace} premise, decoder-success premise, 1024-step generated
unrolling, and monolithic \texttt{inline; sim} are also rejected.

The single Week~10 gate is the actual encoder outer-loop invariant linking
$(i,x,off,encp)$ to the suffix of the pure trace and deriving the success-size
bounds. Only after it compiles should the actual decoder-consumption invariant
be completed. The $h$ codec and security widening remain out of scope.
